\section{考察}
%----------------------------------------------------------------------

\subsection{共有潜在空間の有効性}

本実験の結果から，32次元共有ボトルネック層を介した統合設計が
二重異種性問題に対して有効に機能することが示された．
ボトルネック層を全端末で共有することで，
CNN と ViT という構造的に異なるモデルが同一の潜在表現空間を共有でき，
FedGen の Generator が生成した仮想潜在ベクトルを
端末のアーキテクチャに依らず一様に活用できる．

また，HeteroFL の幅スケーリングをボトルネック層の手前（backbone）のみに適用することで，
ボトルネック以降の重みは全端末で統一され，知識蒸留の学習シグナルが安定する．
この切り分けが，両機構の干渉を防ぐ上で有効に機能していると考えられる．
ただし，ボトルネック設計の有効性をアブレーション実験（ボトルネックなし設定との比較）で
定量的に検証することは今後の課題である．

\subsection{縮小率適応型蒸留の役割}

式(\ref{eq:kd_coeff})に示す蒸留係数の自動調整は，
低能力端末（$r = 0.25$）の精度を補う上で有効な設計と考えられる．
低能力端末は backbone の表現力が制限されているため，
Generator からの知識信号をより強く受け取ることで，
フルサイズ端末との精度ギャップを縮小できると期待される．
なお，FitNet\cite{FitNet}スタイルの中間層蒸留を組み合わせることで，
backbone 自体の表現力を直接補う拡張も今後の方向性として考えられる．

\subsection{ボトルネック次元の選択}

共有ボトルネック層の次元を 32 とした設計上の判断について補足する．
次元を大きくすると，より豊富な潜在表現が共有できる反面，
Generator の訓練に要する計算コストと通信量が増加する．
逆に次元を小さくすると，潜在表現の情報量が不足し知識蒸留の品質が低下するリスクがある．
本研究では 16，32，64 次元での予備実験を行い，
精度と計算コストの均衡点として 32 次元を採用した．
定性的には，16 次元では潜在表現の情報量が不足して精度が低下し，
64 次元では Generator の訓練コストが過大となることを確認した．

また，ボトルネック次元が全端末で固定であるため，
異なるアーキテクチャや異なる幅の端末が同一の潜在空間を共有できる．
これは FedGen の Generator が単一の仮想潜在ベクトル生成器で全端末の蒸留を担えるという，
AFAD のシンプルさの源泉でもある．

\subsection{制約事項}

本稿の評価は IID 設定に限定している．
実環境でより現実的な Non-IID 設定（各端末のデータ分布が偏る状況）では，
生成された仮想潜在ベクトルとローカルデータ分布との乖離が，
知識蒸留の品質に影響を与える可能性がある．
たとえば，局所データに特定クラスが偏っている端末では，
Generator が生成するバランスのとれた仮想潜在ベクトルが
ローカル学習の妨げとなる可能性が考えられる．
さらに，HeteroFL との比較においては参加端末数（5台 vs 10台）および
データ量の差が精度に影響している可能性を否定できない点も留意が必要である．
加えて，Generator の warmup を設けていないため，初期ラウンドでは
Generator が未収束の状態で知識蒸留が行われている可能性があり，
warmup の効果については今後の検証課題とする．
また，FedGen（86.32\%）との精度差 3.4 ポイントについては，
FedGen が計算異種性のない均一モデル設定での評価であるため，
この差が二重異種性に起因するものか設定条件の違いによるものかを切り分けるには，
同条件下での比較検証が別途必要である．
本研究の主張する精度比較はあくまでも HeteroFL（同じ計算異種性あり設定）との比較であり，
FedGen との差については今後の課題とする．

\subsection{他の FL 設定への拡張可能性}

AFAD の共有ボトルネック設計は，CNN と ViT 以外のアーキテクチャ組み合わせ，
たとえば MLP ベースのモデルや Graph Neural Network（GNN）が混在する設定にも
原理的には適用可能である．
ボトルネックを介した潜在空間の共有という設計は backbone の構造に依存しない汎用的な枠組みであり，
今後 FL 環境におけるモデルの多様化が進む中で適用範囲はさらに広がると考えられる．
ただし，各アーキテクチャへの適用に際しては，
bottleneck 次元の最適値や知識蒸留の品質について個別の検証が必要である．

%----------------------------------------------------------------------
